\subsection*{Unit 2 Advanced Math Practice Test}

\setcounter{questionnum}{0}

\noindent\markforreview
\vspace{1em}

For what value of $x$ is the given expression equivalent to $(41k)^{25x}$, where $k > 1$?

\[
\sqrt[8]{41k}\left(\sqrt[9]{41k}\right)^{12}
\]

\vspace{2em}

\noindent\markforreview
\vspace{1em}

The expression $4^{18x}$ is equivalent to $k^{3x}$, where $k$ is a constant. What is the value of $k$?

\vspace{2em}

\noindent\markforreview
\vspace{1em}

The expression $\dfrac{(6x + 8)(3x + 6)}{12x + 24}$ is equivalent to $\dfrac{ax + b}{2}$, where $a$ and $b$ are constants and $x > 0$. What is the value of $a + b$?

\vspace{2em}

\noindent\markforreview
\vspace{1em}

Which expression is equivalent to $\left(3\sqrt{a} - \sqrt{b}\right)^{\frac{2}{7}}$, where $a > b$ and $b > 0$?

\vspace{0.75em}

\choicebox{A}{$\sqrt[7]{9a - b}$}
\choicebox{B}{$\sqrt[7]{9a - 6ab + b}$}
\choicebox{C}{$\sqrt[7]{9a - 3\sqrt{ab} + b}$}
\choicebox{D}{$\sqrt[7]{9a - 6\sqrt{ab} + b}$}

\vspace{2em}

\noindent\markforreview
\vspace{1em}

The expression $12x^2 - mx - 20$, where $m$ is a constant, is equivalent to $(kx - 5)(nx + 4)$, where $k$ and $n$ are integer constants. Which of the following must be an integer?

\vspace{0.75em}

\choicebox{A}{$\dfrac{12}{m}$}
\choicebox{B}{$\dfrac{12}{n}$}
\choicebox{C}{$\dfrac{m}{k}$}
\choicebox{D}{$\dfrac{m}{n}$}

\vspace{2em}

\noindent\markforreview
\vspace{1em}

In the given equation, $m$ and $n$ are constants and $n < 0$. If the equation is true for all values of $x$, what is the value of $m + n$?

\[
4x^2 + mx + 9 = (2x + n)^2
\]

\vspace{0.75em}

\choicebox{A}{$-15$}
\choicebox{B}{$-9$}
\choicebox{C}{$-3$}
\choicebox{D}{$12$}

\vspace{2em}

\noindent\markforreview
\vspace{1em}

The graph of $y = f(x)$ has a $y$-intercept at $(0, 12)$ and $x$-intercepts at $(-3, 0)$ and $(2, 0)$. Which of the following equations could define $f$?

\vspace{0.75em}

\choicebox{A}{$f(x) = (x + 3)^2(x - 2)$}
\choicebox{B}{$f(x) = (x + 3)(x - 2)^2$}
\choicebox{C}{$f(x) = (x - 3)^2(x + 2)$}
\choicebox{D}{$f(x) = (x - 3)(x + 2)^2$}

\vspace{2em}

\noindent\markforreview
\vspace{1em}

Function $f$ is defined by $f(x) = a\sqrt{x} + b$, where $a$ and $b$ are constants. In the $xy$-plane, the graph of $y = f(x)$ has an $x$-intercept at $(16, 0)$, and the graph of $y = f(x) - 10$ has an $x$-intercept at $(9, 0)$. What is the value of $b$?

\vspace{2em}

\noindent\markforreview
\vspace{1em}

For the function $f(x) = b\sqrt{a - x} + 7\sqrt{a - x}$, $a$ and $b$ are constants and $f(16) = 0$. If $f(9) > 0$, which of the following could be the value of $a + b$?

\vspace{0.75em}

\choicebox{A}{$-3$}
\choicebox{B}{$7$}
\choicebox{C}{$9$}
\choicebox{D}{$15$}

\vspace{2em}

\noindent\markforreview
\vspace{1em}

What is the sum of the solutions to the given equation?

\[
x^2 + 7x + 19 = 28 - 5x
\]

\vspace{0.75em}

\choicebox{A}{$-12$}
\choicebox{B}{$-9$}
\choicebox{C}{$-7$}
\choicebox{D}{$12$}

\vspace{2em}

\noindent\markforreview
\vspace{1em}

What are the solutions to the given equation?

\[
4x^2 - 9x - 1 = 0
\]

\vspace{0.75em}

\choicebox{A}{$x = \dfrac{-9 \pm \sqrt{97}}{8}$}
\choicebox{B}{$x = \dfrac{-9 \pm \sqrt{65}}{8}$}
\choicebox{C}{$x = \dfrac{9 \pm \sqrt{97}}{8}$}
\choicebox{D}{$x = \dfrac{9 \pm \sqrt{65}}{8}$}

\vspace{2em}

\noindent\markforreview
\vspace{1em}

One solution to the given equation can be written as $\dfrac{\sqrt{k} - 1}{2}$, where $k$ is a constant. What is the value of $k$?

\[
x^2 + x - 5 = 0
\]

\vspace{2em}

\noindent\markforreview
\vspace{1em}

In the given equation, $k$ is a positive constant. What is the product of the solutions to the given equation?

\[
x^2 - 8kx - (k - 4) = 0
\]

\vspace{0.75em}

\choicebox{A}{$4 - k$}
\choicebox{B}{$k - 4$}
\choicebox{C}{$-8k$}
\choicebox{D}{$8k$}

\vspace{2em}

\noindent\markforreview
\vspace{1em}

If $4x + k$ is a factor of $8x^3 - 26x^2 - 24x$, where $k$ is a positive constant, what is the value of $k$?

\vspace{2em}

\noindent\markforreview
\vspace{1em}

In the given equation, $p$ is a positive constant. The solutions to the given equation are integers. Which of the following is (are) a possible value of $p$?

\[
x^2 + px + 18 = 0
\]

\vspace{0.5em}

\begin{itemize}
\item I. 9
\item II. 17
\item III. 19
\end{itemize}

\vspace{0.75em}

\choicebox{A}{I only}
\choicebox{B}{II only}
\choicebox{C}{I and III only}
\choicebox{D}{I, II, and III}

\vspace{2em}

\noindent\markforreview
\vspace{1em}

The function $f(x) = \dfrac{1}{2}(x - 5)^2 + 8$ gives the profit, in hundreds of thousands of dollars, a company makes when $x$ thousand units of its product are produced and sold, where $1 \leq x \leq 9$. Which of the following is the best interpretation of the vertex of the graph of $y = f(x)$ in the $xy$-plane?

\vspace{0.75em}

\choicebox{A}{The company can make a maximum profit of \$500,000.}
\choicebox{B}{The company can make a maximum profit of \$800,000.}
\choicebox{C}{The company makes a profit of \$0 when 5,000 units of its product are produced and sold.}
\choicebox{D}{The company makes a profit of \$0 when 8,000 units of its product are produced and sold.}

\vspace{2em}

\noindent\markforreview
\vspace{1em}

The given function $P(x) = -0.03(x - 20)^2 + 12$ models the total profit, in thousands of dollars, that a charity group expects to earn from selling boxes of cookies for $x$ dollars each, where $0 \leq x \leq 40$. Based on the model, which of the following is the best interpretation of $P(10) = 9$ in this context?

\vspace{0.75em}

\choicebox{A}{If each box of cookies is sold for \$9, the total profit will likely be \$10,000.}
\choicebox{B}{If each box of cookies is sold for \$9, the total profit will likely be \$21,000.}
\choicebox{C}{If each box of cookies is sold for \$10, the total profit will likely be \$9,000.}
\choicebox{D}{If each box of cookies is sold for \$10, the total profit will likely be \$22,000.}

\vspace{2em}

\noindent\markforreview
\vspace{1em}

In the given system of equations, $b$ is a positive constant. The system has exactly one distinct real solution. What is the value of $b$?

\[
\begin{aligned}
y &= 4x \\
y &= 3x^2 + bx + 3
\end{aligned}
\]

\vspace{2em}

\noindent\markforreview
\vspace{1em}

In the given system of equations, $p$ is a constant. The system has exactly one distinct real solution $(x, y)$. What is the value of $y$?

\[
\begin{aligned}
y &= 2x^2 + 15x + 57 \\
y &= -9x + p
\end{aligned}
\]

\vspace{2em}

\noindent\markforreview
\vspace{1em}

The quadratic function $g(x) = ax^2 + 8x + c$ is defined by the given equation, where $a$ and $c$ are constants. The graph of $y = g(x)$ in the $xy$-plane has $x$-intercepts at $(0, 0)$ and $(n, 0)$, where $n$ is a nonzero constant. If the maximum value of $g(x)$ is 25, which of the following must be true?

\vspace{0.5em}

\begin{itemize}
\item I. $a < 0$
\item II. $c = 0$
\item III. $n > 0$
\end{itemize}

\vspace{0.75em}

\choicebox{A}{I only}
\choicebox{B}{II only}
\choicebox{C}{I and II only}
\choicebox{D}{I, II, and III}

\vspace{2em}

\noindent\markforreview
\vspace{1em}

In the given quadratic function, $b$ and $c$ are constants. The graph of $y = f(x)$ in the $xy$-plane does not intersect the $x$-axis. If $f(-7) = f(9)$, which of the following could be the value of $b + c$?

\[
f(x) = 3x^2 + bx + c
\]

\vspace{0.75em}

\choicebox{A}{$-1$}
\choicebox{B}{$-3$}
\choicebox{C}{$-5$}
\choicebox{D}{$-7$}

\vspace{2em}

\noindent\markforreview
\vspace{1em}

The functions $f$ and $g$ are defined by the given equations, where $x \geq 0$. Which function shows a maximum value as a constant or coefficient for $x \geq 0$?

\vspace{0.5em}

I. $f(x) = 33(0.4)^{x + 3}$

II. $g(x) = 33(0.16)(0.4)^{x - 2}$

\vspace{0.75em}

\choicebox{A}{I only}
\choicebox{B}{II only}
\choicebox{C}{I and II}
\choicebox{D}{Neither I nor II}
